\DocumentMetadata{
  pdfversion=2.0,
  pdfstandard=ua-2,
  lang=en-US,
  tagging=on,
  tagging-setup={math/setup=mathml-SE,tabsorder=structure}
}
\documentclass[11pt,letterpaper]{article}
\usepackage[margin=0.68in,includefoot]{geometry}
\usepackage{newcomputermodern}
\usepackage{amsmath,amscd,mathtools}
\usepackage{tikz,adjustbox}
\providecommand{\square}{\mdlgwhtsquare}
\usepackage[hidelinks]{hyperref}
\hypersetup{
  pdftitle={Algebra PhD Exam, May 1994},
  pdfauthor={University of Florida Department of Mathematics},
  pdfsubject={Graduate mathematics examination},
  pdfdisplaydoctitle=true,
  bookmarksopen=true
}
\setcounter{secnumdepth}{0}
\setlength{\parindent}{0pt}
\setlength{\parskip}{0.28em}
\setlength{\emergencystretch}{3em}
\setlength{\leftmargini}{2.15em}
\setlength{\leftmarginii}{2.65em}
\setlength{\leftmarginiii}{3em}
\pagestyle{empty}
\raggedbottom
\allowdisplaybreaks
\NewTaggingSocketPlug{block/list/label}{examproblem}{%
  \tagstructbegin{tag=Lbl}\tagstructend%
  \tagstructbegin{tag=\LBody}%
  \tagstructbegin{tag=\ProblemHeadingTag,title-o={\ProblemHeadingText}}%
  \tagmcbegin{tag=\ProblemHeadingTag,actualtext={\ProblemHeadingText}}%
  #2%
  \tagmcend%
  \tagstructend%
}
\newenvironment{examproblems}[2]{%
  \def\ProblemHeadingTag{#2}%
  \AssignTaggingSocketPlug{block/list/label}{examproblem}%
  \begin{enumerate}%
  \setcounter{enumi}{#1}%
  \renewcommand{\labelenumi}{\textbf{\theenumi.}}%
  \setlength{\topsep}{0.35em}%
  \setlength{\partopsep}{0pt}%
  \setlength{\itemsep}{0.48em}%
  \setlength{\parsep}{0.18em}%
}{\end{enumerate}}
\newenvironment{examsubparts}{%
  \AssignTaggingSocketPlug{block/list/label}{default}%
  \begin{enumerate}%
  \setlength{\topsep}{0.18em}%
  \setlength{\partopsep}{0pt}%
  \setlength{\itemsep}{0.16em}%
  \setlength{\parsep}{0.1em}%
}{\end{enumerate}}
\newenvironment{examinstructions}{%
  \par\begingroup\itshape
}{%
  \par\endgroup\vspace{0.18em}
}
\makeatletter
\renewcommand\subsection{\@startsection{subsection}{2}{0pt}%
  {-1.25ex plus -.25ex minus -.1ex}%
  {0.55ex plus .1ex}%
  {\normalfont\large\bfseries}}
\renewcommand{\@maketitle}{%
  \newpage\null\vskip -1em%
  {\footnotesize\bfseries
    \href{https://www.ufl.edu/}{University of Florida}, \href{https://math.ufl.edu/}{Department of Mathematics}\par}%
  \vskip -0.5em%
  \tagstructbegin{tag=H1,title={Algebra PhD Exam, May 1994}}%
  \tagpdfparaOff%
  \tagmcbegin{tag=H1}%
    {\LARGE\bfseries\href{https://gma.math.ufl.edu/exams/algebra-phd/}{Algebra PhD Exam}, May 1994\par}%
  \tagmcend%
  \tagpdfparaOn%
  \tagstructend%
  \vskip 0.25em%
  \tagmcbegin{artifact}%
    \hrule height 1pt%
  \tagmcend%
  \vskip 1em%
}
\makeatother
\title{Algebra PhD Exam, May 1994}
\author{}
\date{}
\begin{document}
\maketitle
\thispagestyle{empty}
\begin{examinstructions}
Work 7 out of the following 11 exercises.
\end{examinstructions}
\begin{examproblems}{0}{H2}
\def\ProblemHeadingText{Problem 1.}
\item
Suppose that~\(F\) is a free group on the alphabet~\(X\), and that~\(Y\) is a subset of~\(X\). Let~\(H\) be the least normal subgroup of~\(F\) containing~\(Y\). Prove that \(F/H\) is a free group. (Hint: Show it's free on the alphabet \(X\setminus Y\).)
\def\ProblemHeadingText{Problem 2.}
\item
Suppose that~\(G\) is a finite group.
\begin{examsubparts}
\item[\textnormal{(a)}]
Prove that if~\(G\) is nilpotent, and~\(H\) is any proper subgroup, then~\(H\) is a proper subgroup of its normalizer.
\item[\textnormal{(b)}]
Use (a) to prove that if~\(G\) is nilpotent then it is isomorphic to a finite direct product of~\(p\)-groups.
\end{examsubparts}
\def\ProblemHeadingText{Problem 3.}
\item
Suppose that~\(R\) is a principal ideal domain. Prove that any submodule of a free~\(R\)-module is free.
\def\ProblemHeadingText{Problem 4.}
\item
Prove that there are (up to ring isomorphism) only 12 semisimple left Artinian rings of order 1008, of which only two are not commutative.
\def\ProblemHeadingText{Problem 5.}
\item
Suppose that~\(A\) is a commutative ring with identity. Suppose that \(F(m)\) and \(F(n)\) are free modules on~\(m\) and~\(n\) generators, respectively. Prove that if \(F(m)\cong F(n)\), then \(m=n\).
\def\ProblemHeadingText{Problem 6.}
\item
Suppose that~\(R\) is a ring with identity. Prove that
\[
\operatorname{Hom}_{\mathbb Z}\left(B,\prod_{i\in I}G_i\right)=\prod_{i\in I}\operatorname{Hom}_{\mathbb Z}(B,G_i),
\]
 as right~\(R\)-modules, for all left~\(R\)-modules~\(B\) and all abelian groups \(G_i\) (\(i\in I\)). You may use resources from category theory; if so, outline your argument so that it is clear which theorems you are appealing to.
\def\ProblemHeadingText{Problem 7.}
\item
Let~\(A\) be a commutative ring with identity. For each multiplicative system~\(S\) of~\(A\), prove that \(S^{-1}A\) is a flat~\(A\)-module.
\def\ProblemHeadingText{Problem 8.}
\item
State and prove the Hilbert Nullstellensatz. (You may use the following: if~\(K\) is an algebraically closed field, \(A\) a finitely generated~\(K\)-algebra, and~\(M\) is a maximal ideal of~\(A\), then \(A/M\) is isomorphic to~\(K\).)
\def\ProblemHeadingText{Problem 9.}
\item
Among the following integral domains, decide which ones are Dedekind domains, and give a brief explanation. Quoting a relevant theorem will do; likewise, a well-illustrated example.
\begin{examsubparts}
\item[\textnormal{(a)}]
\(\mathbb Z[T]\), the polynomial ring over the integers, in one variable.
\item[\textnormal{(b)}]
\(\mathbb Z[\sqrt{-5}]=\{a+b\sqrt{-5}:a,b\in\mathbb Z\}\).
\item[\textnormal{(c)}]
The ring \(k[[T]]\) of all formal power series in one variable, over the field~\(k\).
\end{examsubparts}
\def\ProblemHeadingText{Problem 10.}
\item
Suppose that~\(E\) is a finite Galois extension of the field~\(F\). If \(\operatorname{Gal}(E/F)\) has order \(pq\), where \(p<q\) are primes, such that~\(p\) does not divide \(q-1\), prove that~\(E\) has two subfields \(E_p\) and \(E_q\), which are stable under the action of \(\operatorname{Gal}(E/F)\), such that \(E_p\cap E_q=F\), \(E_p\) and \(E_q\) generate~\(E\), and \(\operatorname{Gal}(E_p/F)\) (resp. \(\operatorname{Gal}(E_q/F)\)) is cyclic of order~\(p\) (resp. \(q\)).
\def\ProblemHeadingText{Problem 11.}
\item
This problem has two parts.
\begin{examsubparts}
\item[\textnormal{(a)}]
Define: algebraic closure.
\item[\textnormal{(b)}]
Prove that every field has an algebraic closure, and argue that if~\(E\) is an algebraic closure of~\(F\), then \(|E|\) is countable if~\(F\) is finite, while \(|E|=|F|\), otherwise.
\end{examsubparts}
\end{examproblems}
\end{document}
